% =====================================================================
\section{The analogy with ACTN (RFC 8453)}
\label{sec:actn}
% =====================================================================

GS~QKD~018 does not invent a new model: it applies to QKD a pattern that is
already standard in transport networks, the IETF ACTN framework
(\emph{Abstraction and Control of TE Networks})~\cite{rfc8453}. ACTN has
three levels: the \emph{Customer Network Controller} (CNC), the
\emph{Multi-Domain Service Coordinator} (MDSC) and the per-domain
\emph{Provisioning Network Controllers} (PNC), connected by the CMI
(CNC–MDSC) and MPI (MDSC–PNC) interfaces. The mapping to GS~QKD~018 is
direct (Table~\ref{tab:actn}).

\begin{table}[h]
\centering\small
\caption{Mapping between ACTN and the ETSI QKD stack.}
\label{tab:actn}
\begin{tabularx}{\textwidth}{@{}l l L@{}}
\toprule
\textbf{ACTN} & \textbf{ETSI QKD} & \textbf{Remark} \\
\midrule
CNC & administrator / customer portal & originates the service request \\
MDSC & SDN orchestrator & multi-domain view, path decision \\
PNC & SDN controller of the QKD domain & computes candidates, activates paths \\
MPI & GS~QKD~018 interface & abstract topology, provisioning, notifications \\
SBI & GS~QKD~015 interface & control of individual nodes \\
abstract TE topology & \texttt{etsi-qkd-sdn-topology} & vendor-neutral view \\
TE bandwidth & SKR / ESKR & the resource to allocate is key \\
\bottomrule
\end{tabularx}
\end{table}

The essential difference lies in the resource. In ACTN the bandwidth of a
link is fixed by the channel capacity; in QKD the SKR fluctuates with
physical conditions and key accumulates in finite buffers from which it is
consumed. This is why in 018 notifications such as \texttt{\_perf\_update}
and \texttt{\_overloaded} and the \texttt{path\_restoration} policy play a
more central role than in classical transport networks.
